\documentclass[../main.tex]{subfiles}

\begin{document}
% -----------------------------
\section{Conclusion}
% -----------------------------
    This tutorial demonstrated the complete journey of compiling a simple machine learning
computation from a custom dialect to an executable program. We began by defining a HiCompiler dialect,
introducing domain-specific operations that represent high-level computations.
Next, we implemented lowering rules that translate these operations into standard MLIR dialects
such as arith, enabling the program to move closer to hardware-level representation.
After lowering, we applied compiler optimizations like constant folding and dead code
elimination to simplify the program and improve efficiency. Once the IR was normalized
and optimized, we used MLIR passes to convert the program into the LLVM dialect,
which bridges MLIR and the LLVM compiler infrastructure. The LLVM dialect was then
translated into LLVM IR, a low-level, platform-independent intermediate representation.
    
    Using LLVM tools, we compiled the LLVM IR into assembly and object files, preparing
it for execution on the target machine. We then linked the compiled object with
a small C harness, allowing the generated function to be called from a standard executable
program. This final step enabled us to run and test the compiled model, confirming that
the entire compilation pipeline works end to end. Overall, the tutorial illustrated how
modern compiler infrastructures like MLIR and LLVM allow developers to design domain-specific
languages, perform powerful optimizations, and ultimately generate efficient native executables.

\end{document}
